% Versao: v2 (13/11/2026)
\documentclass[12pt]{article}

\usepackage{sbc-template}
\usepackage{graphicx}
\usepackage[hyphens]{url}
\usepackage[utf8]{inputenc}
\usepackage[brazilian]{babel}
\usepackage{listings}
\usepackage{amssymb}
\usepackage{tabularx}
\usepackage{tikz}
\usetikzlibrary{positioning,arrows.meta}
\usepackage[portuguese,onelanguage,ruled,vlined,nokwfunc]{algorithm2e}

% Nota: o modelo original traz \usepackage[latin1]{inputenc} apos a chamada com
% utf8. Como a segunda chamada sobrepoe a primeira e quebraria a acentuacao,
% mantem-se apenas utf8.

\lstset{
  basicstyle=\ttfamily\footnotesize,
  breaklines=true,
  breakatwhitespace=false,
  columns=fullflexible,
  keepspaces=true,
  frame=none,
  aboveskip=8pt,
  belowskip=8pt,
  xleftmargin=0.2cm
}

% algorithm2e: palavras-chave em portugues (portuguese + onelanguage), sem
% ponto e virgula no fim das linhas, fonte dois pontos menor que o texto e
% comentarios em italico, para que cada bloco caiba numa pagina (um bloco [H]
% nao quebra entre paginas).
\SetAlFnt{\footnotesize}
\DontPrintSemicolon
\SetCommentSty{textit}
\SetKwFunction{LimparEstado}{LimparEstado}
\SetKwFunction{AmostraPorComprimento}{AmostraPorComprimento}
\SetKwFunction{RelatarSimilaridades}{RelatarSimilaridades}

% Identificadores e caminhos longos em fonte de maquina de escrever, com quebra
% de linha permitida depois de hifen, barra e sublinhado.
\DeclareUrlCommand\caminho{\urlstyle{tt}}

% Coluna de largura automatica do tabularx, alinhada a esquerda, para que o
% texto das celulas nao fique com espacos esticados.
\newcolumntype{Y}{>{\raggedright\arraybackslash}X}

\sloppy

\title{Predizendo a taxa SELIC a partir de resumos de atas do COPOM\\
  \normalsize Relatório parcial de Iniciação Científica\\
  \normalsize Versão v2 (13/11/2026)}

\author{Giovanni Willik Del Piccolo\inst{1}}

\address{Escola de Artes, Ciências e Humanidades (EACH)\\
  Universidade de São Paulo (USP) -- São Paulo -- SP -- Brasil
  \email{giovanniwillik@usp.br}
}

\begin{document}

\maketitle

\begin{abstract}
  This partial report describes an undergraduate research project that asks
  whether summaries of the COPOM minutes, written by a language model running
  locally, predict the direction of the SELIC rate better than the minutes
  truncated by \textit{head+tail}, the input of the reference work, which
  fine-tuned BERT and BERTimbau. The hypothesis is that a summary fits the
  input limit of these classifiers and keeps more information than truncation
  does. Initial studies fixed the summarising model and the prompt and showed
  that, under greedy decoding and a uniform server state, repeated submissions
  of a minute yield the same summary, so one submission per minute suffices.
  Algorithm 1 produced the corpus of summaries in Portuguese and English. The
  classification of the summaries and the comparison with the reference work
  have not been executed yet.
\end{abstract}

\begin{resumo}
  Este relatório parcial descreve a Iniciação Científica que investiga se
  resumos das atas do COPOM, escritos por um modelo de língua executado
  localmente, predizem a direção da taxa SELIC melhor do que as atas truncadas
  por \textit{head+tail}, entrada do trabalho de referência, que fez o ajuste
  fino de BERT e BERTimbau. A hipótese é que o resumo cabe no limite de entrada
  desses classificadores e preserva mais informação do que o truncamento.
  Estudos iniciais fixaram o modelo sumarizador e o prompt e mostraram que, com
  decodificação gulosa e estado do servidor uniforme, submissões repetidas de
  uma ata produzem o mesmo resumo, de modo que basta uma submissão por ata. O
  Algoritmo 1 gerou o corpus de resumos em português e em inglês. A
  classificação dos resumos e a comparação com o trabalho de referência ainda
  não foram executadas.
\end{resumo}

\section{Introdução}\label{sec:intro}

O Comitê de Política Monetária (COPOM) do Banco Central do Brasil publica uma
ata de cada reunião. A pergunta que motiva esta Iniciação Científica é se o
texto de uma ata ajuda a prever a direção da taxa SELIC na reunião seguinte:
alta, manutenção ou queda.

O ponto de partida é o Trabalho de Conclusão de Curso de Forti~\cite{forti2025},
que fez o ajuste fino de BERT~\cite{devlin2019}, para as atas em inglês, e de
BERTimbau~\cite{souza2020}, para as atas em português, com 30 repetições de
\textit{holdout}. Como esses modelos aceitam no máximo 512 tokens e as atas são
muito mais longas, cada ata foi truncada por \textit{head+tail}: ficam os 255
primeiros e os 255 últimos tokens, e o meio é descartado, estratégia proposta
por Sun et al.~\cite{sun2019}. Os resultados foram F1-macro ${\approx}$ 0,25 e
acurácia ${\approx}$ 40\% nos dois idiomas, e o modelo falha na classe
manutenção.

A hipótese desta pesquisa é que um resumo da ata inteira, escrito por um modelo
de língua de grande escala (LLM) executado localmente, cabe nos mesmos 512
tokens e preserva mais informação do que o truncamento, porque o meio da ata,
que o \textit{head+tail} descarta, traz o balanço de riscos e as indicações do
Comitê sobre os próximos passos. O objetivo é verificar se esses resumos são
preditores melhores do que as atas truncadas. O projeto impõe três condições: o
LLM reside na máquina local e não muda durante o estudo; o prompt é definido de
antemão, sem escolher entre prompts pelo resultado da classificação (as
restrições R1 a R3 da seção~\ref{sec:prompt}); e a avaliação usa a mesma
metodologia de Forti~\cite{forti2025}.

Não foram encontrados trabalhos que usem um LLM local para resumir atas de
banco central antes da classificação com BERT, nem que comparem, em documentos
de política monetária, a sumarização por LLM com o truncamento
\textit{head+tail} (seção~\ref{sec:relacionados}). Esta pesquisa combina a
sumarização como pré-processamento~\cite{mutasodirin2024} com a previsão de
política monetária~\cite{shah2023,fasolo2022} e a aplica às atas do COPOM.

Depois dos trabalhos relacionados (seção~\ref{sec:relacionados}) e da
metodologia geral (seção~\ref{sec:metodologia}), a seção~\ref{sec:estudos}
reúne os estudos que antecederam os algoritmos, e a
seção~\ref{sec:sumarizacao} especifica a sumarização e relata a sua execução.
As seções~\ref{sec:classificacao} e~\ref{sec:resultados-gerais} ficam em branco
nesta versão, porque a classificação ainda não foi executada. A
seção~\ref{sec:limitacoes} reúne as limitações e os próximos passos.

\section{Trabalhos relacionados}\label{sec:relacionados}

O processamento de linguagem natural aplicado à comunicação de bancos
centrais estuda sobretudo o Federal Reserve e o Banco Central Europeu. Na
classificação do tom (\textit{hawkish}, \textit{dovish} ou neutro),
Shah et al.~\cite{shah2023} mostraram que modelos da família BERT captam esses
sinais em sentenças do FOMC, embora a tarefa continue difícil. Na previsão de
decisões, Jingyi e Liu~\cite{jingyi2025} combinaram representações TF-IDF de
textos do FOMC com indicadores econômicos e chegaram a AUC de 0,83, acima dos
modelos que usam só uma das fontes. Sobre o Banco Central do Brasil, Fasolo et
al.~\cite{fasolo2022} mostraram que o tom das atas e dos comunicados do COPOM
tem efeito mensurável sobre expectativas de inflação, atividade e incerteza. O
trabalho mais próximo desta pesquisa é o de Forti~\cite{forti2025}, que é o seu
\textit{baseline}.

O BERT aceita 512 tokens. Sun et al.~\cite{sun2019} compararam formas de
truncar documentos longos e concluíram que manter o início e o fim
(\textit{head+tail}) é a melhor delas, base da escolha de Forti. Mutasodirin e
Prasojo~\cite{mutasodirin2024} compararam truncamento e sumarização como formas
de encurtar documentos antes de classificá-los com BERT: os resumos superaram a
maioria das variantes de truncamento, perderam só para manter o início do
documento e, em alguns conjuntos de dados, ficaram próximos do desempenho com o
texto completo. Os sumarizadores usados ali não eram LLMs. Uma alternativa sem
sumarização é dividir o documento em trechos, processar cada um com BERT e
combinar as saídas~\cite{pappagari2019}, ao custo de uma arquitetura mais
complexa.

Entre os LLMs abertos, o LLaMA-3-8B teve resultados competitivos na
sumarização, sobretudo de textos longos~\cite{rehman2025}. O risco principal é a
fidelidade: Kim et al.~\cite{kim2024} encontraram erros de fidelidade em mais da
metade dos resumos de livros produzidos pelos LLMs avaliados. As atas são muito
mais curtas que livros, o que deve atenuar o problema, mas não o elimina. Em
português, os modelos abertos maiores têm desempenho razoável, abaixo do obtido
em inglês, com destaque para as famílias Llama 3 e Qwen~\cite{almeida2025}.

\section{Metodologia geral}\label{sec:metodologia}

O trabalho encadeia três estágios (figura~\ref{fig:pipeline}). A sumarização
resume cada ata com um LLM local; a classificação faz o ajuste fino de BERT e
BERTimbau sobre os resumos; a verificação compara o desempenho obtido com o do
\textit{baseline} de Forti~\cite{forti2025}, que usou as atas truncadas. O
modelo sumarizador, o prompt e o número de submissões por ata foram fixados
antes, pelos estudos da seção~\ref{sec:estudos}.

\begin{figure}[ht]
\centering
\begin{tikzpicture}[
    node distance=4mm,
    caixa/.style={draw, rounded corners=2pt, align=center, font=\small,
      text width=23mm, minimum height=15mm, inner sep=2pt},
    futura/.style={caixa, dashed},
    dado/.style={align=center, font=\small, text width=23mm, inner sep=2pt},
    seta/.style={-{Stealth[length=2mm]}, semithick}]
  \node[dado] (atas) {Atas do COPOM\\em PT e EN\\(202 cada)};
  \node[caixa, right=of atas] (sum) {Sumarização\\Algoritmo 1\\Llama 3.1 8B};
  \node[dado, right=of sum] (res) {Tabela\\\texttt{Resumo}\\(404 resumos)};
  \node[futura, right=of res] (cla) {Classificação\\Algoritmo 2\\BERT e\\BERTimbau};
  \node[futura, right=of cla] (ver) {Verificação\\resumo contra\\truncamento};
  \node[caixa, above=7mm of sum, text width=46mm, minimum height=10mm]
    (est) {Estudos iniciais (seção~\ref{sec:estudos}):\\modelo, prompt e
    \texttt{N\_sub} = 1};
  \node[dado, below=7mm of ver, text width=38mm] (base)
    {\textit{Baseline} de Forti:\\\textit{head+tail}, 30 repetições};
  \draw[seta] (atas) -- (sum);
  \draw[seta] (sum) -- (res);
  \draw[seta] (res) -- (cla);
  \draw[seta] (cla) -- (ver);
  \draw[seta] (est) -- (sum);
  \draw[seta] (base) -- (ver);
\end{tikzpicture}
\caption{Visão geral do trabalho. As caixas tracejadas são os estágios ainda
não executados.}
\label{fig:pipeline}
\end{figure}

\paragraph{Sumarização.} O Algoritmo 1 (seção~\ref{sec:sumarizacao}) recebe as
atas dos dois idiomas e devolve a tabela \texttt{Resumo}, com o texto literal
de um resumo por ata, e a lista \texttt{Acima} dos resumos que passam de 512
tokens no tokenizador do classificador de destino. Os resumos não são editados.

\paragraph{Classificação.} O Algoritmo 2 (seção~\ref{sec:classificacao}) recebe
a tabela \texttt{Resumo} como dado de entrada e pode ser executado sem executar
o Algoritmo 1: os dois algoritmos são independentes, por exigência do
orientador. Ele replica o protocolo de Forti, com 30 repetições de
\textit{holdout}, divisão 80/20 estratificada e feita por reunião, semente 42,
três recortes temporais e hiperparâmetros fixos, e muda apenas a entrada: o
resumo em vez da ata truncada. Na ingestão, remove-se o preâmbulo dos resumos em
inglês e corta-se o final dos que passam de 512 tokens, pelas regras decididas
com o orientador (seção~\ref{sec:decisoes}).

\paragraph{Verificação.} Com a mesma semente e o mesmo número de repetições, a
repetição $i$ desta pesquisa usa a mesma partição da repetição $i$ do
\textit{baseline}, cujos resultados estão no repositório de
Forti~\cite{forti2025}. A comparação é feita em 30 pares (resumo, truncamento)
por idioma e recorte, e também por classe, em particular na classe manutenção,
em que o \textit{baseline} falha.

\subsection{Dados}\label{sec:dados}

O corpus é o de Forti~\cite{forti2025}: 202 atas do COPOM em cada idioma, de
janeiro de 2002 a março de 2025. As atas em português foram obtidas pela API do
Banco Central do Brasil, e as em inglês, pelos PDFs do site do Banco Central.
O pré-processamento de Forti normaliza o texto pela forma NFKD, remove os
caracteres que não são ASCII e colapsa os espaços. Cada ata é o texto completo e
limpo de uma reunião, com a sua data de referência. Os dois conjuntos são as
entradas \texttt{D\_pt} e \texttt{D\_en} dos procedimentos das
seções~\ref{sec:determinismo} e~\ref{sec:alg1}.

O alvo da classificação é a direção da decisão da reunião seguinte, em
\{alta, manutenção, queda\}, construída a partir das datas e das decisões das
reuniões. O estudo usa três recortes temporais, com corte em 2016-01-01:
\texttt{before\_2016}, com 129 atas por idioma, \texttt{after\_2016}, com 73, e
\texttt{full}, com todas. O conjunto de atas usado para escolher o prompt é
outro, posterior a março de 2025 e disjunto deste (seção~\ref{sec:prompt}).

\subsection{Ambiente e reprodução}\label{sec:ambiente}

Todas as execuções usaram o ambiente da tabela~\ref{tab:ambiente}. O código,
os resumos e as saídas estão no repositório público do
projeto~\cite{delpiccolo2026}, na pasta \caminho{ic_exploracao/}; as atas de
entrada e os resultados do \textit{baseline}, no repositório de
Forti~\cite{forti2025}.

\begin{table}[ht]
\centering
\caption{Ambiente de execução do estudo}
\label{tab:ambiente}
\begin{tabularx}{\linewidth}{|l|Y|}
\hline
\textbf{Componente} & \textbf{Versão} \\
\hline
Hardware & MacBook Air M4, 24 GB RAM, macOS 26.6.2 arm64 \\
\hline
LLM & \texttt{llama3.1:8b}, Q4\_K\_M, digest 46e0c10c039e --- o mesmo em todas
as fases \\
\hline
Servidor & Ollama 0.32.15 na escolha do prompt; 0.34.2 no piloto do experimento
de determinismo (05/10/2026); 0.35.1 na primeira execução do experimento e na
replicação controlada (12/10/2026); 0.40.5 na segunda execução (19/10/2026);
no corpus, ver a seção~\ref{sec:desvios} \\
\hline
Python & 3.10.0 \\
\hline
transformers / torch & 4.44.2 / 2.4.1 \\
\hline
sentence-transformers & 3.0.1, modelo
\caminho{paraphrase-multilingual-MiniLM-L12-v2} \\
\hline
Tokenizadores & \caminho{bert-base-portuguese-cased} (PT),
\caminho{bert-base-cased} (EN) \\
\hline
\end{tabularx}
\end{table}

\section{Estudos iniciais}\label{sec:estudos}

Três decisões foram tomadas antes de construir o Algoritmo 1, e ele as recebe
prontas: o modelo sumarizador (seção~\ref{sec:escolha-llm}), o prompt
(seção~\ref{sec:prompt}) e o número de submissões por ata
(seção~\ref{sec:determinismo}). As duas últimas valem para a combinação de
servidor e modelo da tabela~\ref{tab:ambiente}.

\subsection{Escolha do modelo sumarizador}\label{sec:escolha-llm}

O LLM precisa residir na máquina local, resumir bem em português e em inglês por
um prompt fixado de antemão, gerar resumos que caibam nos 512 tokens do BERT e
caber num MacBook Air M4 com 24 GB de memória. Foram considerados o Llama 3.1
8B, o Qwen 2.5 (7 e 14 bilhões de parâmetros), o Mistral 7B, treinado sobretudo
em inglês, e o Gemma 7B, descartado porque o seu contexto de 8 mil tokens é
menor que o das atas mais longas do corpus (seção~\ref{sec:num_ctx}).

O modelo escolhido foi o Llama 3.1 8B Instruct~\cite{grattafiori2024}. A
sumarização de documentos longos por modelos dessa família tem avaliação
publicada~\cite{rehman2025}; o contexto de 128 mil tokens recebe
qualquer ata inteira, sem truncar a entrada do sumarizador; o desempenho em
português está entre os melhores dos modelos abertos~\cite{almeida2025}; a
licença (Meta Llama 3.1 Community License) permite o uso acadêmico; e a versão
quantizada Q4\_K\_M ocupa cerca de 4,7 GB e deixa folga de memória na máquina.
A estabilidade entre execuções foi verificada depois, no experimento da
seção~\ref{sec:determinismo}. O Qwen 2.5 7B ficou como alternativa e não
precisou ser usado.

\subsection{Escolha do prompt}\label{sec:prompt}

O projeto de pesquisa impõe três restrições ao prompt:

\begin{itemize}
\item R1 --- O prompt deve ser definido a priori, em português e em inglês.
\item R2 --- É vedado testar vários prompts e selecionar os que favorecem o
  resultado final da tarefa (falácia da evidência incompleta).
\item R3 --- O resumo deve caber nos 512 tokens de entrada do BERT/BERTimbau.
\end{itemize}

R1 e R2 combinadas implicam que a seleção não pode observar métricas da
classificação: todo critério precisa ser intrínseco ao resumo. Usar atas do
recorte do estudo contaminaria o experimento, porque as mesmas atas seriam
depois usadas no treino e na avaliação do classificador. Por isso foram
coletadas as 11 atas posteriores a março de 2025 (reuniões 270 a 280), pela API
pública do Banco Central e com o pré-processamento de Forti. A avaliação usou
três delas --- 270 (20.955 caracteres), 274 (15.166) e 279 (16.650) ---, que
cobrem os extremos e o meio da faixa de comprimento.

\paragraph{Candidatos e critérios.} Quatro candidatos foram escritos antes de
qualquer execução (tabela~\ref{tab:candidatos}); nenhum pede a predição da
decisão futura ou um julgamento de postura. Os textos integrais estão no
repositório do projeto~\cite{delpiccolo2026}, em
\caminho{ic_exploracao/validacao_prompt/testar_prompts.py}. Os critérios,
todos intrínsecos ao resumo, foram quatro: V1, orçamento de tokens (no máximo
512 no tokenizador de destino, restrição rígida); V2, neutralidade (nenhum
termo que antecipe a decisão seguinte, por lista definida a priori); V3,
fidelidade (nenhuma informação ausente do documento, por leitura); e V4,
economia do orçamento (proporção dos 512 tokens ocupada por texto constante).

\begin{table}[ht]
\centering
\caption{Candidatos de prompt e respectivas estratégias de instrução}
\label{tab:candidatos}
\begin{tabularx}{\linewidth}{|l|Y|}
\hline
\textbf{Candidato} & \textbf{Estratégia de instrução} \\
\hline
A --- Neutro & Pede apenas o resumo, com limite de extensão e cláusula de
fidelidade. Não nomeia temas. \\
\hline
B --- Temático & Enumera os temas a cobrir quando presentes (externo,
atividade, trabalho, inflação, crédito, riscos). \\
\hline
C --- Estruturado & Impõe gabarito de cinco seções fixas, três frases por
seção, com marcação ``Não abordado''. \\
\hline
D --- Extrativo & Pede reprodução literal dos dez trechos mais informativos,
sem reformulação. \\
\hline
\end{tabularx}
\end{table}

\paragraph{Avaliação.} Os candidatos foram submetidos ao \texttt{llama3.1:8b},
no servidor da fase de seleção (tabela~\ref{tab:ambiente}), com
\texttt{temperature = 0}, \texttt{top\_p = 1}, \texttt{num\_predict = 900},
\texttt{num\_ctx = 16384} e limite de 350 palavras
(tabela~\ref{tab:prompt_resultados}). O \texttt{num\_ctx} desta rodada trunca
parte das atas do corpus do estudo (seção~\ref{sec:num_ctx}), mas não as de
validação, que têm no máximo 20.955 caracteres, cerca de 6.200 tokens.

\begin{table}[ht]
\centering
\caption{Resultados medidos nas atas 270, 274 e 279}
\label{tab:prompt_resultados}
\begin{tabular}{|l|l|l|l|l|}
\hline
\textbf{Candidato} & \textbf{Tokens BERT} & \textbf{Cabe em 512} &
\textbf{Cobertura} & \textbf{Viés} \\
 & \textbf{(270 / 274 / 279)} & & \textbf{temática} & \\
\hline
A --- Neutro & 358 / 258 / 338 & sim (3/3) & 5 / 5 / 6 de 7 & nenhum \\
\hline
B --- Temático & 243 / 284 / 287 & sim (3/3) & 5 / 6 / 6 de 7 & nenhum \\
\hline
C --- Estruturado & 355 / 417 / 352 & sim (3/3) & 7 / 7 / 7 de 7 & nenhum \\
\hline
D --- Extrativo & 506 / 427 / 552 & não (2/3) & 6 / 6 / 5 de 7 & nenhum \\
\hline
\end{tabular}
\end{table}

\paragraph{Decisão.} D foi eliminado por V1: produziu 552 tokens na ata 279. A,
B e C satisfizeram V1, V2 e V3 nas três atas, e C foi selecionado de início pela
cobertura temática de 7/7. A escolha foi revertida por três motivos: o critério
de cobertura é circular, porque a lista dos sete temas foi definida pelo
pesquisador e C enumera esses temas no próprio prompt; nomear temas, como fazem
B e C, injeta uma hipótese do pesquisador sobre o que é relevante, e uma
eventual melhora sobre o \textit{baseline} deixaria de ser atribuível à
sumarização; e o gabarito de C gasta o orçamento com texto constante (cerca de
22 tokens de cabeçalhos por resumo, mais 12 a 24 marcadores, contra V4), com uma
marcação ``Não abordado'' que poderia aparecer de forma desigual antes e depois
de 2016.
\textbf{Prompt selecionado: candidato A (neutro).} Ele isola a variável em
estudo, gera de 258 a 358 tokens e produz prosa contínua, sem cabeçalhos. A
sumarização guiada por domínio, com C, fica como extensão futura
(seção~\ref{sec:limitacoes}).

\paragraph{Limite de palavras.} A tokenização WordPiece fragmenta palavras, e a
razão entre tokens e palavras medida nos resumos gerados foi, em média, 1,37 em
português e 1,42 em inglês, com pior caso de 1,48 numa ata inglesa de 2013. Dos
512 tokens, 510 são úteis, e o limite seguro é 510 / 1,48 $\approx$ 344
palavras. O limite foi reduzido de 350 para 300 palavras: 300 $\times$ 1,48 = 444
tokens, com 66 de folga. A garantia de R3 vem da medição de cada resumo no
Algoritmo 1 e do corte na ingestão (seção~\ref{sec:decisoes}).

\paragraph{Prompt fixado.} Os prompts abaixo são imutáveis para o restante da
pesquisa; a implementação confere o SHA-256 do prompt e para se ele mudar. O
marcador \texttt{\{m\}} é substituído pelo texto completo e limpo da ata, nunca
por um fragmento. A versão inglesa é tradução literal da portuguesa.

\noindent\textbf{Português} (\texttt{Prompt\_pt})
\begin{lstlisting}
Resuma o documento a seguir em no maximo 300 palavras.
O resumo deve preservar os temas discutidos e as avaliacoes apresentadas, sem
acrescentar informacao que nao esteja no texto original.

Documento:
{m}

Resumo:
\end{lstlisting}

\noindent\textbf{Inglês} (\texttt{Prompt\_en})
\begin{lstlisting}
Summarize the document below in at most 300 words.
The summary must preserve the topics discussed and the assessments presented,
without adding information that is not in the original text.

Document:
{m}

Summary:
\end{lstlisting}

\subsection{Experimento de determinismo}\label{sec:determinismo}

Este experimento, feito uma vez antes de gerar o corpus, fixou o número de
submissões por ata usado pelo Algoritmo 1. Não é parte do algoritmo nem é
repetido a cada execução dele.

\subsubsection{Motivação e critério de decisão}\label{sec:det-motivacao}

O projeto previa submeter cada ata várias vezes ao mesmo prompt para medir a
variação das respostas, e a primeira versão do Algoritmo 1 fixava, por isso,
\texttt{N\_sub = 10} submissões por ata. Por determinação do orientador, esse
número passou a ser decidido por evidência, antes de gerar o corpus, com
critério definido antes da execução: submeter dez atas dez vezes cada, nos dois
idiomas, e observar se os resumos de uma mesma ata saem idênticos, caractere
por caractere. Saindo idênticos, a repetição não carrega informação e passa-se
a uma submissão por ata; saindo diferentes, mantêm-se dez. Medir a
variabilidade com \texttt{temperature} $> 0$ foi descartado pelo orientador,
porque aumenta o risco de o modelo acrescentar informação ausente da ata.

\subsubsection{Mecanismo}\label{sec:mecanismo}

\paragraph{Decodificação gulosa e semente.} Em cada passo, a rede devolve um
\textit{logit} por token do vocabulário, e um procedimento posterior, o
amostrador, escolhe o próximo token~\cite{holtzman2020}. Com
\texttt{temperature} $> 0$, o próximo token é sorteado, e a semente determina o
sorteio. Com \texttt{temperature} $= 0$, o amostrador toma sempre o token de
maior logit (decodificação gulosa): nenhum número pseudoaleatório é consumido,
e \textbf{a semente nunca é lida}. Fixar a semente, portanto, não torna uma
submissão reproduzível; o que a torna reproduzível é o estado de execução do
servidor.

\paragraph{Ponto flutuante e cache de prefixo.} A adição em ponto flutuante não
é associativa~\cite{goldberg1991}, e a ordem das somas que produzem cada logit
depende de quantos tokens são processados de uma vez. O servidor Ollama guarda
o processamento do último prompt (o cache de prefixo). Quando as submissões de
uma ata são feitas em sequência, a primeira encontra no cache o prompt da ata
anterior e processa o texto do zero; as seguintes reaproveitam o prompt
idêntico. São dois caminhos de execução, com logits que diferem nos últimos
bits. Quando os dois primeiros candidatos estão quase empatados, essa diferença
inverte a escolha do token, e os dois textos seguem caminhos distintos dali em
diante. Limpar o estado do servidor antes de cada chamada faz todas as chamadas
seguirem o mesmo caminho.

\subsubsection{Procedimento do teste}\label{sec:det-procedimento}

O teste chama as sub-rotinas abaixo: \texttt{LimparEstado}, que o Algoritmo 1
também usa, a amostragem das atas e as três similaridades da tabela
\texttt{Metricas}, que recebem dois resumos da mesma ata e valem 1 para textos
idênticos. A implementação está em
\caminho{ic_exploracao/algoritmo1/determinismo.py} e
\caminho{ic_exploracao/algoritmo1/analisar_variancia.py}, no repositório do
projeto~\cite{delpiccolo2026}.

\begin{itemize}
\item \texttt{LimparEstado}: pede ao servidor que descarregue o modelo (no
  Ollama, uma requisição com \texttt{keep\_alive = 0} e sem prompt), consulta a
  lista de modelos residentes (\texttt{GET /api/ps}) até ela ficar vazia, por
  até 30 s, e interrompe a execução se o modelo continuar residente, porque
  seguir misturaria dois estados do servidor. Ao retornar, o cache de prefixo
  está vazio.
\item \texttt{AmostraPorComprimento(D, K)}: ordena as atas de \texttt{D} pelo
  número de caracteres, com empate desfeito pela data, e devolve as datas das
  posições $\mathrm{round}(i \cdot (n - 1)/(K - 1))$, para
  $i = 0, \ldots, K - 1$, contadas a partir de 0, sendo $n$ o número de atas. É
  aplicada a \texttt{D\_pt}, e as mesmas reuniões são usadas em inglês.
\item \texttt{SimilaridadeLexical} (métrica \texttt{lexical}): índice de
  Jaccard~\cite{jaccard1912} entre os conjuntos de palavras dos dois resumos,
  depois de dobrar o texto para ASCII (forma NFKD), passá-lo para minúsculas e
  trocar por espaço o que não é letra ou dígito; vale 0 se a união for vazia.
\item \texttt{SimilaridadeRougeL} (métrica \texttt{rougeL}): ROUGE-L na
  variante F1~\cite{lin2004}, a média harmônica da precisão e da revocação da
  maior subsequência comum entre as sequências de palavras, pela mesma
  transformação.
\item \texttt{SimilaridadeEmbedding} (métrica \texttt{embedding}): cosseno
  entre os embeddings de sentença dos dois textos literais, pela média das
  saídas de \texttt{Emb}~\cite{reimers2019}, que é fixo e nunca é retreinado.
\item \texttt{RelatarSimilaridades(Sim, lang)}: para cada métrica, o mínimo, os
  três quartis (interpolação linear), o máximo, a amplitude interquartil, a
  média e o desvio padrão (\texttt{ddof = 1}), no geral, por ata, por ano e por
  recorte temporal.
\end{itemize}

O teste é o procedimento~\ref{proc:teste}. Ele recebe as atas dos dois idiomas,
os prompts fixados, o modelo sumarizador com os seus parâmetros, o tamanho da
amostra \texttt{K}, o número de submissões por ata \texttt{N\_test}, a semente
base \texttt{S0} e o passo \texttt{P} entre sementes, o modelo de embeddings
\texttt{Emb}, a tabela \texttt{Metricas}, que associa cada nome de métrica à
função que a calcula (lista acima), e o número \texttt{N\_sub\_max}
de submissões por ata que o corpus teria se o teste encontrasse variação. Ele
devolve a decisão \texttt{N\_sub}, o número de resumos distintos de cada ata da
amostra e as similaridades entre as submissões, com as suas estatísticas. Os
resumos do teste são gravados em arquivo próprio e nunca entram no corpus. Os
passos são:
% PROSA proc:teste
\begin{enumerate}
\item \texttt{Identicas} recebe verdadeiro; ela indica se todas as submissões
  saíram iguais.
\item \texttt{Amostra} recebe as datas das \texttt{K} reuniões escolhidas por
  \texttt{AmostraPorComprimento}, que as espalha pela ordem de comprimento das
  atas em português.
\item Para cada idioma \texttt{lang}.
\item Para cada ata \texttt{m} da amostra, tomada de \texttt{D\_pt} ou
  \texttt{D\_en} conforme o idioma.
\item Para cada submissão \texttt{j} de 1 até \texttt{N\_test}.
\item \texttt{LimparEstado} descarrega o modelo e confirma que ele saiu da
  memória do servidor.
\item \texttt{s\_j} recebe \texttt{S0} + \texttt{P}$\cdot$\texttt{j}: é a
  semente da submissão \texttt{j}.
\item \texttt{T[lang][m][j]} recebe o resumo do \texttt{LLM} para a ata
  \texttt{m}, em sessão nova e com \texttt{keep\_alive = 0}: \texttt{T} guarda
  os resumos do teste.
\item \texttt{Distintos[lang][m]} recebe quantos textos diferentes as
  submissões de \texttt{m} produziram; se mais de um, \texttt{Identicas} recebe
  falso.
\item Para cada par \texttt{(a, b)} de submissões e cada métrica de
  \texttt{Metricas}, \texttt{Sim[nome][lang][m][a][b]} recebe a similaridade
  entre \texttt{a} e \texttt{b}.
\item \texttt{RelatarSimilaridades} calcula as estatísticas de cada métrica no
  idioma e as grava no relatório do teste.
\item \texttt{N\_sub} recebe 1 se \texttt{Identicas} for verdadeira e
  \texttt{N\_sub\_max} se não: é o número de submissões por ata do corpus.
\end{enumerate}

\begin{procedure}[H]
\caption{TesteDeterminismo()}\label{proc:teste}
\KwIn{\texttt{Idiomas} = \{pt, en\}; \texttt{D\_pt}, \texttt{D\_en}: atas limpas
  (seção~\ref{sec:dados})\\
  \texttt{Prompt\_pt}, \texttt{Prompt\_en}: prompts fixados
  (seção~\ref{sec:prompt}); o marcador \texttt{\{m\}} é substituído
  pelo texto completo e limpo da ata \texttt{m}, nunca por um fragmento\\
  \texttt{LLM}: \texttt{llama3.1:8b}, Q4\_K\_M, digest 46e0c10c039e, no Ollama;
  \texttt{temperature = 0}, \texttt{top\_p = 1}, \texttt{num\_predict = 900},
  \texttt{num\_ctx = 32768}\\
  \texttt{K} = 10: atas da amostra; \texttt{N\_test} = 10: submissões por ata\\
  \texttt{S0} = 42 e \texttt{P} = 9973: semente base e passo entre sementes\\
  \texttt{Emb}: \texttt{paraphrase-multilingual-MiniLM-L12-v2}\\
  \texttt{Metricas}: \texttt{lexical}, \texttt{rougeL} e \texttt{embedding}
  (seção~\ref{sec:det-procedimento})\\
  \texttt{N\_sub\_max} = 10: submissões por ata se houver variação}
\KwOut{\texttt{N\_sub}; \texttt{Identicas}; \texttt{Distintos[lang][m]};
  \texttt{Sim[nome][lang][m][a][b]} e as suas estatísticas}
\tcp{\texttt{Identicas} recebe verdadeiro; ela indica se todas as submissões
  saíram iguais.}
\texttt{Identicas} $\gets$ verdadeiro\;
\tcp{\texttt{Amostra} recebe as datas das \texttt{K} reuniões escolhidas por
  \texttt{AmostraPorComprimento}, que as espalha pela ordem de comprimento das
  atas em português.}
\texttt{Amostra} $\gets$ \AmostraPorComprimento{\texttt{D\_pt}, \texttt{K}}\;
\tcp{Para cada idioma \texttt{lang}.}
\ForEach{idioma \texttt{lang} de \texttt{Idiomas}}{
  \tcp{Para cada ata \texttt{m} da amostra, tomada de \texttt{D\_pt} ou
    \texttt{D\_en} conforme o idioma.}
  \ForEach{ata \texttt{m} de \texttt{D\_lang} com data em \texttt{Amostra}}{
    \tcp{Para cada submissão \texttt{j} de 1 até \texttt{N\_test}.}
    \For{$\texttt{j} \gets 1$ \KwTo \texttt{N\_test}}{
      \tcp{\texttt{LimparEstado} descarrega o modelo e confirma que ele saiu da
        memória do servidor.}
      \LimparEstado{}\;
      \tcp{\texttt{s\_j} recebe \texttt{S0} + \texttt{P}$\cdot$\texttt{j}: é a
        semente da submissão \texttt{j}.}
      \texttt{s\_j} $\gets$ \texttt{S0} + \texttt{P}$\cdot$\texttt{j}\;
      \tcp{\texttt{T[lang][m][j]} recebe o resumo do \texttt{LLM} para a ata
        \texttt{m}, em sessão nova e com \texttt{keep\_alive = 0}: \texttt{T}
        guarda os resumos do teste.}
      \texttt{T[lang][m][j]} $\gets$ \texttt{LLM}(\texttt{Prompt\_lang} com
      \texttt{m}, semente \texttt{s\_j}, \texttt{keep\_alive = 0})\;
    }
    \tcp{\texttt{Distintos[lang][m]} recebe quantos textos diferentes as
      submissões de \texttt{m} produziram; se mais de um, \texttt{Identicas}
      recebe falso.}
    \texttt{Distintos[lang][m]} $\gets$ número de textos distintos entre
    \texttt{T[lang][m][1..N\_test]}\;
    \lIf{\texttt{Distintos[lang][m]} $> 1$}{\texttt{Identicas} $\gets$ falso}
    \tcp{Para cada par \texttt{(a, b)} de submissões e cada métrica de
      \texttt{Metricas}, \texttt{Sim[nome][lang][m][a][b]} recebe a similaridade
      entre \texttt{a} e \texttt{b}.}
    \ForEach{par \texttt{(a, b)} com $1 \le a < b \le$ \texttt{N\_test} e
      \texttt{(nome, Similaridade)} de \texttt{Metricas}}{
      \texttt{Sim[nome][lang][m][a][b]} $\gets$
      \texttt{Similaridade}(\texttt{T[lang][m][a]}, \texttt{T[lang][m][b]})\;
    }
  }
  \tcp{\texttt{RelatarSimilaridades} calcula as estatísticas de cada métrica no
    idioma e as grava no relatório do teste.}
  \RelatarSimilaridades{\texttt{Sim}, \texttt{lang}}\;
}
\tcp{\texttt{N\_sub} recebe 1 se \texttt{Identicas} for verdadeira e
  \texttt{N\_sub\_max} se não: é o número de submissões por ata do corpus.}
\leIf{\texttt{Identicas}}{\texttt{N\_sub} $\gets$ 1}{\texttt{N\_sub} $\gets$
  \texttt{N\_sub\_max}}
\end{procedure}

\subsubsection{Execuções e resultados}\label{sec:det-resultados}

O procedimento~\ref{proc:teste} foi executado duas vezes sobre a mesma amostra,
com os mesmos prompts, parâmetros e sementes (tabela~\ref{tab:det_execucoes}).
A amostra cobre atas de 5.807 a 67.477 caracteres e, com o \texttt{num\_ctx}
adaptativo da seção~\ref{sec:desvios}, os três patamares de \texttt{num\_ctx} em
português e dois em inglês. A primeira execução foi feita sem o passo 6
(limpar o estado): as dez submissões de cada ata em sequência, com o modelo
residente. Foi ela que revelou a dependência do estado do cache e levou a
incluir a limpeza no procedimento. A segunda execução seguiu o procedimento como
especificado.

\begin{table}[ht]
\centering
\caption{As duas execuções do teste de determinismo}
\label{tab:det_execucoes}
\begin{tabularx}{\linewidth}{|>{\raggedright\arraybackslash}p{3.4cm}|Y|Y|}
\hline
 & \textbf{Primeira execução (sem limpar o estado)} &
\textbf{Segunda execução (estado limpo)} \\
\hline
Estado do servidor & modelo residente; cache de prefixo reaproveitado & limpo
antes de cada chamada \\
\hline
Data e servidor & 12/10/2026, Ollama 0.35.1 & 19/10/2026, Ollama 0.40.5 \\
\hline
Chamadas & 200, sem falhas & 200, sem falhas \\
\hline
Combinações idênticas & 17 de 20 (PT 9 de 10, EN 8 de 10) & \textbf{20 de 20} \\
\hline
Combinações que variaram & pt 20230920, en 20070606, en 20110720 & nenhuma \\
\hline
Tempo de processamento & 2,5 h & 7,9 h (5,1 h PT, 2,7 h EN) \\
\hline
Tempo por chamada & --- & mediana 128 s, mínimo 28 s, máximo 330 s \\
\hline
\end{tabularx}
\end{table}

\paragraph{Primeira execução e replicação controlada.} Nas três combinações que
variaram, a submissão $j = 1$ diferiu, e as submissões $j = 2$ a $j = 10$ foram
idênticas entre si. Esse padrão já é evidência contra a semente: as nove
submissões idênticas usam nove sementes diferentes. Para isolar a causa, a
submissão $j = 1$ de cada caso foi refeita com a mesma semente, duas vezes
seguidas: com o modelo recém-carregado, a chamada reproduziu $j = 1$; com o
prompt idêntico já no cache, reproduziu $j = 2$ a $j = 10$, nos três casos. O
estado do cache determina a saída; a semente, não. A diferença vai de uma frase
reescrita (pt 20230920, Jaccard de 0,9925) a um resumo substancialmente outro
(en 20110720, Jaccard de 0,3818). Os textos divergentes estão no repositório do
projeto, em \caminho{ic_exploracao/resumos/determinismo_padrao_en.jsonl} e no
arquivo equivalente em português.

\paragraph{Segunda execução.} Com o estado limpo antes de cada chamada, as dez
submissões de cada combinação produziram um único resumo, idêntico caractere
por caractere, nas vinte combinações. Jaccard e ROUGE-L valem exatamente 1,0
nos 450 pares de cada idioma; o cosseno difere de 1 só na sétima casa decimal,
por arredondamento da precisão simples (\texttt{float32}). Nas 20 combinações, o
resumo em estado limpo é idêntico à submissão $j = 1$ da primeira execução: a
primeira submissão de cada ata segue o caminho de uma chamada em estado limpo,
e a variante divergente era a das submissões que reaproveitavam o cache. Com
\texttt{temperature = 0} e estado de execução uniforme, o sumarizador é uma
função determinística da ata.

\subsubsection{Equivalência entre a primeira chamada e o estado
limpo}\label{sec:equivalencia}

O Ollama se atualizou sozinho durante o estudo (tabela~\ref{tab:ambiente}). A
segunda execução foi feita inteira na versão 0.40.5, estabelecida pelo horário
de partida do processo do servidor (instalação em 16/10/2026, às 12h27, sem
reinício até o fim da execução), e o resumo em estado limpo nessa versão é
idêntico, nas 20 combinações, ao da submissão $j = 1$ gerada nas versões 0.34.2
e 0.35.1. Na 0.40.5, a equivalência entre a primeira chamada de uma ata e o
estado limpo foi verificada nas três combinações sensíveis, chamando antes a ata
que precedeu a alvo na geração do corpus: a primeira chamada reproduziu o estado
limpo nos três casos, e uma segunda chamada reproduziu a variante que
reaproveita o cache observada na 0.35.1. Como cada ata do corpus é submetida
uma única vez, encontrando no cache o prompt de outra ata, o resumo que o corpus
registra é o que uma chamada em estado limpo produz. Isso só deixaria de valer
se a geração fosse interrompida no meio de uma ata e retomada; por isso o
Algoritmo 1 limpa o estado na partida e antes de cada nova tentativa (passos 1
e 9).

\subsubsection{Decisão e limitação}\label{sec:det-decisao}

\textbf{Decisão:} \textbf{\texttt{N\_sub = 1}}, uma submissão por ata, saída do
procedimento~\ref{proc:teste} na segunda execução. A repetição não mede
variabilidade do modelo: a semente é inerte com \texttt{temperature = 0}, a
diferença da primeira execução é artefato do cache do servidor e, removido o
artefato, a variância é zero. A decisão foi estabelecida uma vez, para a
combinação de servidor e modelo da tabela~\ref{tab:ambiente}; o Algoritmo 1 a
recebe como entrada e não testa o determinismo. Com uma submissão por ata, a
análise de variância sobre o corpus, prevista na primeira versão do algoritmo,
perdeu o objeto e não foi executada.

\textbf{Limitação:} a equivalência foi verificada nas 20 combinações da amostra,
e não nas 202 atas de cada idioma, e depende do servidor, da versão do Ollama e
do modelo (\texttt{llama3.1:8b}, digest 46e0c10c039e). Se algum deles mudar, a
verificação precisa ser refeita (\texttt{determinismo.py --verificar-versao});
por isso a versão do servidor e o digest do modelo são gravados em cada resumo.

\section{Sumarização automática}\label{sec:sumarizacao}

\subsection{Metodologia}\label{sec:alg1}

O Algoritmo 1 resume cada ata do corpus, nos dois idiomas, com o modelo
sumarizador e o prompt fixados na seção~\ref{sec:estudos}, e mede o
comprimento de cada resumo no tokenizador do classificador de destino. Esta
subseção diz como reproduzi-lo; a seção~\ref{sec:res-alg1} relata o que a execução
produziu. O algoritmo é independente do Algoritmo 2, que recebe a sua saída como
dado.

O algoritmo recebe as atas dos dois idiomas (seção~\ref{sec:dados}), os prompts
fixados (seção~\ref{sec:prompt}), o modelo sumarizador com os seus
parâmetros, o número de submissões por ata \texttt{N\_sub}, o limite de entrada
do classificador \texttt{L\_max} e as sementes. \texttt{N\_sub} vale 1: é a
conclusão do experimento de determinismo (seção~\ref{sec:det-decisao}), válida
para o servidor e o modelo da tabela~\ref{tab:ambiente}. O algoritmo continua
escrito para qualquer valor de \texttt{N\_sub}. Ele devolve a tabela
\texttt{Resumo}, com o texto literal de cada resumo, e, para cada idioma, a
lista \texttt{Acima} dos resumos que passam de \texttt{L\_max} tokens. A sub-rotina
\texttt{LimparEstado} está definida na seção~\ref{sec:det-procedimento}. Os passos são:

% PROSA alg:sumarizacao
\begin{enumerate}
\item \texttt{LimparEstado} descarrega o modelo e confirma que ele saiu da
  memória, para que a geração parta sem nenhum prompt no cache.
\item Para cada idioma \texttt{lang}.
\item \texttt{Acima[lang]} começa vazia; ela vai listar os resumos do idioma
  com mais de \texttt{L\_max} tokens.
\item Para cada ata \texttt{m} de \texttt{D\_pt} ou de \texttt{D\_en},
  conforme o idioma.
\item Para cada submissão \texttt{j} de 1 até \texttt{N\_sub}.
\item Se \texttt{N\_sub} $> 1$, \texttt{LimparEstado} descarrega o modelo, para
  que todas as submissões da ata processem o prompt do zero.
\item Conferir a versão do servidor e parar a geração se ela mudou desde a
  partida.
\item \texttt{s\_j} recebe \texttt{S0} + \texttt{P}$\cdot$\texttt{j}: é a
  semente da submissão \texttt{j}.
\item \texttt{Resumo[lang][m][j]} recebe o resumo do \texttt{LLM} para a ata
  \texttt{m}, em sessão nova; se a chamada falhar, \texttt{LimparEstado} vem
  antes de cada nova tentativa.
\item \texttt{n} recebe o número de tokens do resumo, com \texttt{[CLS]} e
  \texttt{[SEP]}, no tokenizador do BERTimbau
  (\texttt{bert-base-portuguese-cased}) se pt e do BERT
  (\texttt{bert-base-cased}) se en.
\item Se \texttt{n} passa de \texttt{L\_max}, acrescentar \texttt{(m, j, n)} a
  \texttt{Acima[lang]}; o resumo fica como saiu, e o corte é feito pelo
  Algoritmo 2.
\item Relatar quantos resumos há em \texttt{Acima[lang]} e a distribuição do
  excesso \texttt{n} $-$ \texttt{L\_max}.
\end{enumerate}

\begin{algorithm}[H]
\caption{Sumarização das atas}\label{alg:sumarizacao}
\KwIn{\texttt{Idiomas} = \{pt, en\}; \texttt{D\_pt}, \texttt{D\_en}: atas limpas
  (seção~\ref{sec:dados})\\
  \texttt{Prompt\_pt}, \texttt{Prompt\_en}: prompts fixados
  (seção~\ref{sec:prompt}); o marcador \texttt{\{m\}} é substituído
  pelo texto completo e limpo da ata \texttt{m}, nunca por um fragmento\\
  \texttt{LLM}: \texttt{llama3.1:8b}~\cite{grattafiori2024}, Q4\_K\_M, digest
  46e0c10c039e, no Ollama; \texttt{temperature = 0}, \texttt{top\_p = 1},
  \texttt{num\_predict = 900}, \texttt{num\_ctx = 32768}
  (seção~\ref{sec:num_ctx})\\
  \texttt{N\_sub} = 1: submissões por ata, fixado pelo experimento de
  determinismo (seção~\ref{sec:det-decisao})\\
  \texttt{L\_max} = 512: limite de entrada do BERT~\cite{devlin2019}, com
  \texttt{[CLS]} e \texttt{[SEP]}\\
  \texttt{S0} = 42 e \texttt{P} = 9973: semente base e passo entre sementes}
\KwOut{\texttt{Resumo[lang][m][j]}: texto literal de cada resumo, gravado com a
  versão do servidor e o digest do modelo\\
  \texttt{Acima[lang]}: resumos com mais de \texttt{L\_max} tokens}
\tcp{\texttt{LimparEstado} descarrega o modelo e confirma que ele saiu da
  memória, para que a geração parta sem nenhum prompt no cache.}
\LimparEstado{}\;
\tcp{Para cada idioma \texttt{lang}.}
\ForEach{idioma \texttt{lang} de \texttt{Idiomas}}{
  \tcp{\texttt{Acima[lang]} começa vazia; ela vai listar os resumos do idioma
    com mais de \texttt{L\_max} tokens.}
  \texttt{Acima[lang]} $\gets$ lista vazia\;
  \tcp{Para cada ata \texttt{m} de \texttt{D\_pt} ou de \texttt{D\_en},
    conforme o idioma.}
  \ForEach{ata \texttt{m} de \texttt{D\_lang}}{
    \tcp{Para cada submissão \texttt{j} de 1 até \texttt{N\_sub}.}
    \For{$\texttt{j} \gets 1$ \KwTo \texttt{N\_sub}}{
      \tcp{Se \texttt{N\_sub} $> 1$, \texttt{LimparEstado} descarrega o modelo,
        para que todas as submissões da ata processem o prompt do zero.}
      \lIf{\texttt{N\_sub} $> 1$}{\LimparEstado{}}
      \tcp{Conferir a versão do servidor e parar a geração se ela mudou desde
        a partida.}
      \lIf{a versão do servidor mudou desde a partida}{parar a geração}
      \tcp{\texttt{s\_j} recebe \texttt{S0} + \texttt{P}$\cdot$\texttt{j}: é a
        semente da submissão \texttt{j}.}
      \texttt{s\_j} $\gets$ \texttt{S0} + \texttt{P}$\cdot$\texttt{j}\;
      \tcp{\texttt{Resumo[lang][m][j]} recebe o resumo do \texttt{LLM} para a
        ata \texttt{m}, em sessão nova; se a chamada falhar,
        \texttt{LimparEstado} vem antes de cada nova tentativa.}
      \texttt{Resumo[lang][m][j]} $\gets$ \texttt{LLM}(\texttt{Prompt\_lang}
      com \texttt{m}, semente \texttt{s\_j})\;
      \tcp{\texttt{n} recebe o número de tokens do resumo, com \texttt{[CLS]} e
        \texttt{[SEP]}, no tokenizador do BERTimbau
        (\texttt{bert-base-portuguese-cased}) se pt e do BERT
        (\texttt{bert-base-cased}) se en.}
      \texttt{n} $\gets$ número de tokens de \texttt{Resumo[lang][m][j]}\;
      \tcp{Se \texttt{n} passa de \texttt{L\_max}, acrescentar
        \texttt{(m, j, n)} a \texttt{Acima[lang]}; o resumo fica como saiu, e o
        corte é feito pelo Algoritmo 2.}
      \lIf{\texttt{n} $>$ \texttt{L\_max}}{acrescentar \texttt{(m, j, n)} a
        \texttt{Acima[lang]}}
    }
  }
  \tcp{Relatar quantos resumos há em \texttt{Acima[lang]} e a distribuição do
    excesso \texttt{n} $-$ \texttt{L\_max}.}
  relatar \texttt{Acima[lang]}\;
}
\end{algorithm}

\paragraph{Notas do Algoritmo 1.} As notas trazem só justificativa; o
algoritmo pode ser reproduzido sem elas.

{\renewcommand{\theenumi}{\roman{enumi}}
\renewcommand{\labelenumi}{(\theenumi)}
\begin{enumerate}
\item \textbf{Sementes.} As sementes são espaçadas por um passo \texttt{P} primo,
  precaução contra fluxos pseudoaleatórios correlacionados~\cite{lecuyer2002}
  que só teria efeito com \texttt{temperature} $> 0$
  (seção~\ref{sec:mecanismo}).
\item \textbf{Resumos acima de \texttt{L\_max}.} Regerar exigiria mudar o
  prompt ou a semente, ambos fixos, e reamostrar só os casos acima do limite
  enviesaria a distribuição de comprimento para resumos mais curtos. O limite é
  imposto na ingestão pelo Algoritmo 2 (seção~\ref{sec:decisoes}), e a lista
  \texttt{Acima} mostra a frequência e a magnitude dos cortes.
\item \textbf{Limpeza do estado.} Com \texttt{N\_sub = 1}, cada submissão
  encontra no cache o prompt de outra ata, nunca o da própria, caminho que
  equivale ao estado limpo (seção~\ref{sec:equivalencia}). A limpeza na partida
  e antes de cada nova tentativa evita que uma ata interrompida deixe parte do
  próprio prompt no cache, e a versão do servidor é conferida porque a
  equivalência vale para uma combinação específica de servidor e modelo.
\end{enumerate}}

A implementação está em \caminho{ic_exploracao/algoritmo1/gerar_resumos.py},
no repositório do projeto~\cite{delpiccolo2026}, e os comandos de execução, em
\caminho{ic_exploracao/algoritmo1/README.md}.

\subsection{Resultados}\label{sec:res-alg1}

As contagens desta subseção foram recalculadas a partir dos arquivos de dados do
repositório do projeto~\cite{delpiccolo2026}.

\subsubsection{Execução e integridade}\label{sec:corpus-execucao}

O algoritmo foi executado sobre o corpus inteiro (tabela~\ref{tab:corpus_execucao}),
sem nenhuma ata faltando ou duplicada e com o mesmo prompt (conferido pelo
SHA-256) e o mesmo modelo (conferido pelo digest) em todos os resumos.

\begin{table}[ht]
\centering
\caption{Execução do Algoritmo 1}
\label{tab:corpus_execucao}
\begin{tabularx}{\linewidth}{|Y|c|c|}
\hline
 & \textbf{PT} & \textbf{EN} \\
\hline
Resumos & 202 & 202 \\
\hline
Falhas & 0 & 0 \\
\hline
Atas truncadas na entrada do LLM & 0 & 0 \\
\hline
Atas por patamar de \texttt{num\_ctx} (8k / 16k / 32k) & 70 / 65 / 67 & 102 / 100 / 0 \\
\hline
Elevações de patamar por truncamento detectado & 0 & 0 \\
\hline
Tempo por resumo (mediana / máximo) & 144 s / 756 s & 157 s / 537 s \\
\hline
Tempo de processamento & 10,3 h & 9,4 h \\
\hline
\end{tabularx}
\end{table}

\subsubsection{Comprimento dos resumos}\label{sec:comprimento}

O limite de 300 palavras não é respeitado com precisão
(tabela~\ref{tab:comprimento}): o modelo passa dele em 27,2\% dos resumos em
português e 46,0\% em inglês. A garantia de comprimento é a medição em tokens
do Algoritmo 1 (passos 10 e 11). Os resumos em português trazem acentuação (201
de 202), embora as atas de entrada sejam ASCII: o modelo restaura os acentos ao
escrever, o que justifica o dobramento para ASCII nas similaridades léxica e
ROUGE-L (seção~\ref{sec:det-procedimento}).

\begin{table}[ht]
\centering
\caption{Comprimento dos resumos (Q1 e Q3: primeiro e terceiro quartis)}
\label{tab:comprimento}
\small
\begin{tabularx}{\linewidth}{|Y|c|c|}
\hline
 & \textbf{PT} & \textbf{EN} \\
\hline
Tokens BERT: mín. / Q1 / mediana / Q3 / máx. &
201 / 311 / 362 / 414 / 808 & 249 / 351 / 389 / 467 / 912 \\
\hline
Tokens BERT, média & 378 & 421 \\
\hline
Palavras, mediana & 268 & 296 \\
\hline
Acima das 300 palavras pedidas & 55 (27,2\%) & 93 (46,0\%) \\
\hline
Compressão em caracteres, mediana & 5,3\% & 6,8\% \\
\hline
\end{tabularx}
\end{table}

\subsubsection{Resumos acima de \texttt{L\_max}}\label{sec:acima}

O inglês tem cerca de 2,6 vezes mais resumos acima do limite que o português
(tabela~\ref{tab:over}). Todos os 51 casos são de atas anteriores a 2016: as
atas mais antigas são mais longas e produzem resumos mais longos (mediana de 403
contra 305 tokens em português, e de 428 contra 353 em inglês). O corte na
ingestão pelo Algoritmo 2 incide, portanto, apenas sobre o recorte
\texttt{before\_2016}, e com intensidade diferente em cada idioma.

\begin{table}[ht]
\centering
\caption{Resumos que excedem 512 tokens (lista \texttt{Acima}; Q1 e Q3: primeiro e
terceiro quartis)}
\label{tab:over}
\small
\begin{tabularx}{\linewidth}{|Y|c|c|}
\hline
 & \textbf{PT} & \textbf{EN} \\
\hline
Acima de 512 tokens & 14 de 202 (6,9\%) & 37 de 202 (18,3\%) \\
\hline
Excesso: mín. / Q1 / mediana / Q3 / máx. &
7 / 54 / 62 / 135 / 296 & 2 / 28 / 66 / 101 / 400 \\
\hline
Excesso, média & 94 & 97 \\
\hline
Antes de 2016 & 14 de 129 (10,9\%) & 37 de 129 (28,7\%) \\
\hline
A partir de 2016 & 0 de 73 & 0 de 73 \\
\hline
\end{tabularx}
\end{table}

\subsubsection{O que o corte do final remove}\label{sec:corte}

O corte do final mantém os primeiros 510 tokens do resumo, mais \texttt{[CLS]}
e \texttt{[SEP]} (seção~\ref{sec:decisoes}). Nos 51 resumos afetados:

\begin{itemize}
\item o corte remove, na mediana, 11\% do resumo, com máximo de 37\% em
  português e 44\% em inglês;
\item o trecho removido contém termos da decisão de política monetária (Selic,
  decisão, manter, elevar, reduzir e equivalentes em inglês) em 4 dos 14
  resumos em português e 20 dos 37 em inglês --- mas, em todos os 51, esses
  termos também aparecem na parte mantida;
\item a alternativa \textit{head+tail} do \textit{baseline} (255 tokens
  iniciais e 255 finais) removeria o meio do resumo, que contém esses termos em
  6 dos 14 e 31 dos 37.
\end{itemize}

A medição é por palavras-chave: indica onde a decisão é mencionada, mas não
garante que o trecho removido seja redundante.

\subsubsection{Saídas degeneradas}\label{sec:degeneradas}

Em duas atas, pt 20030423 (808 tokens BERT, 477 palavras) e en 20110608 (912
tokens, 582 palavras), o modelo passou a reproduzir trechos do documento em vez
de resumi-lo e só parou no limite \texttt{num\_predict = 900}. No caso
português a cópia é visível: o final do resumo está sem acentos, como a ata de
entrada. Pelo passo 11 do Algoritmo 1, as duas saídas ficam como produzidas e
são cortadas na ingestão.

\subsubsection{Preâmbulo nos resumos em inglês}\label{sec:preambulo}

55 dos 202 resumos em inglês (27,2\%) começam por uma frase fixa, sem conteúdo
da ata, em uma de duas formas. Cada forma é uma única linha, com um espaço
simples entre as palavras; a quebra abaixo é só de diagramação. As duas
coincidem até ``less'' e diferem no caractere seguinte: vírgula na longa,
dois-pontos na curta.

\begin{itemize}
\item \texttt{Pre\_long}, em 18 resumos, com 22 tokens BERT: \texttt{Here is a
  summary of the document in 300 words or less, preserving the topics discussed
  and the assessments presented:}
\item \texttt{Pre\_short}, em 37 resumos, com 13 tokens BERT: \texttt{Here is a
  summary of the document in 300 words or less:}
\end{itemize}

Nos 55 casos a frase abre o texto e é seguida de duas quebras de linha; ela não
aparece em outra posição, e nenhum resumo em português a tem. Dos 55, 51 são de
atas anteriores a 2016 (o primeiro é de 18/07/2007 e o último de 08/06/2016). A
regra de remoção (seção~\ref{sec:decisoes}), aplicada aos dois arquivos do
corpus sem gravá-los, remove exatamente 55 preâmbulos em inglês (18 da forma
longa e 37 da curta) e nenhum em português. Depois da remoção, a mediana dos
resumos em inglês passa de 389 para 378 tokens BERT, e a média, de 421 para 417;
o máximo continua em 912. Nenhum dos 37 resumos em inglês acima do limite tem
preâmbulo: a remoção não muda quais resumos são cortados nem quanto se corta de
cada um.

\subsubsection{Desvios em relação à especificação}\label{sec:desvios}

\textbf{Versões do servidor.} Por causa da atualização automática do Ollama
(seção~\ref{sec:equivalencia}), o corpus mistura três versões: em português, 14
resumos em 0.34.2, 8 em 0.35.1 e 180 em 0.40.5; em inglês, 10 em 0.35.1 e 192
em 0.40.5, sempre com o mesmo modelo. Os 32 resumos das versões antigas incluem
os 20 da amostra do experimento de determinismo, idênticos aos da 0.40.5; os
outros 12, todos em português, não têm essa contraprova. Desde 20/10/2026 a
geração confere a versão antes de cada chamada e para se ela mudar (passo 7).
\textbf{\texttt{num\_ctx} adaptativo.} Em vez do valor único da entrada
\texttt{LLM}, cada ata usou o menor patamar de \{8192, 16384, 32768\} que a
acomodasse sem truncar, no corpus e no experimento de determinismo, depois de
medido que o resumo sai idêntico nos três patamares; nenhuma ata exigiu
elevação de patamar (tabela~\ref{tab:corpus_execucao}).

\subsubsection{Decisões tomadas com o orientador}\label{sec:decisoes}

Três pontos foram enviados ao orientador em 20/10/2026, cada um com uma
proposta. As três propostas foram aprovadas, e as respostas foram registradas em
02/11/2026. As duas últimas definem a ingestão dos resumos pelo Algoritmo 2.

\paragraph{Correção do \texttt{num\_ctx}.}\label{sec:num_ctx} A primeira versão
da entrada \texttt{LLM} do Algoritmo 1 declarava \texttt{num\_ctx = 16384}. O
Ollama não recusa um prompt maior que \texttt{num\_ctx}: processa apenas
\texttt{num\_ctx / 2} tokens e descarta o resto, sem aviso. Na maior ata
portuguesa do corpus (67.477 caracteres), o campo \texttt{prompt\_eval\_count}
mostrou 4.098 tokens processados com \texttt{num\_ctx = 8192}, 8.194 com
\texttt{16384} e os 19.945 do documento com \texttt{32768}. Com
\texttt{num\_ctx = 16384}, 49 das 202 atas portuguesas (24,3\%) seriam
truncadas e nenhuma das inglesas, porque o português gera mais tokens por
caractere (3,38 caracteres por token, contra 4,57 em inglês), o que confundiria
a comparação entre idiomas. Com \texttt{32768}, nenhuma ata é truncada, e o
modelo usa cerca de 9~GB de memória. \textbf{Proposta:} manter a correção.
\textbf{Resposta do orientador:} aprovada (``OK''). O valor está na entrada
\texttt{LLM} do Algoritmo~\ref{alg:sumarizacao}.

\paragraph{Corte do final dos resumos longos.} \textbf{Proposta:} cortar o
final dos resumos acima de 512 tokens, em vez de usar o \textit{head+tail} do
\textit{baseline}, porque o corte do final remove pouco e preserva os termos da
decisão, que o \textit{head+tail} tiraria com mais frequência
(seção~\ref{sec:corte}). \textbf{Resposta do orientador:} aprovada (``OK''); o
\textit{head+tail} fica descartado. \textbf{Regra:} se o resumo tokenizado, com
\texttt{[CLS]} e \texttt{[SEP]}, passa de 512 tokens, mantêm-se \texttt{[CLS]},
os 510 primeiros tokens do resumo e \texttt{[SEP]}, e o corte é registrado.

\paragraph{Remoção do preâmbulo.} \textbf{Proposta:} remover o preâmbulo dos
resumos em inglês na ingestão, por regra fixa, mantendo o corpus literal
(seção~\ref{sec:preambulo}). \textbf{Resposta do orientador:} remover
(``Remove''). \textbf{Regra:} se o resumo em inglês começa exatamente por
\texttt{Pre\_long}, remove-se \texttt{Pre\_long}; senão, se começa exatamente
por \texttt{Pre\_short}, remove-se \texttt{Pre\_short}; nos dois casos,
removem-se também os espaços e as quebras de linha seguintes, e nada mais. A
comparação é de cadeia exata e só no início do texto. Como nenhuma forma é
prefixo da outra, a ordem de teste não muda o resultado; a longa é testada
primeiro para que a regra tenha uma única leitura.

\paragraph{Ordem da ingestão.} Os passos têm ordem fixa: (1) remover o
preâmbulo, só em inglês; (2) tokenizar com o tokenizador do idioma, o do
BERTimbau (\texttt{bert-base-portuguese-cased}) em português e o do BERT
(\texttt{bert-base-cased}) em inglês, com \texttt{[CLS]} e \texttt{[SEP]}; (3)
cortar o final até 512 tokens. O corte vem depois da remoção, para que o
preâmbulo não ocupe o orçamento; como a remoção só encurta o texto, todo resumo
cortado está na lista \texttt{Acima}. As regras agem sobre a cópia carregada
pelo Algoritmo 2: a tabela \texttt{Resumo} não é alterada.

\section{Classificação}\label{sec:classificacao}

\subsection{Metodologia}\label{sec:alg2}

\subsection{Resultados}\label{sec:res-alg2}

\section{Resultados gerais}\label{sec:resultados-gerais}

\section{Limitações e próximos passos}\label{sec:limitacoes}

\paragraph{Limitações.} O prompt foi escolhido em três atas, e o resumo inglês
não foi avaliado em dados disjuntos, porque a API do Banco Central não expõe as
atas em inglês posteriores a março de 2025; a fidelidade, sujeita a erros em
LLMs~\cite{kim2024}, foi verificada só por leitura de amostra. A equivalência
que dispensa limpar o estado foi verificada numa amostra e vale para uma
combinação de servidor e modelo (seção~\ref{sec:det-decisao}), e o corpus
mistura três versões do servidor (seção~\ref{sec:desvios}). Os cortes na
ingestão se concentram em \texttt{before\_2016} e em inglês
(seção~\ref{sec:acima}), o que pesa na leitura dos resultados por recorte e por
idioma. Por fim, 202 atas por idioma é pouco para o ajuste fino de um modelo com
110 milhões de parâmetros: a sumarização melhora a entrada, não o volume de
treino. Um resultado sem melhoria sobre o \textit{baseline} também é uma
contribuição válida e será relatado como tal.

\paragraph{Próximos passos.} Na ordem em que serão feitos:

\begin{enumerate}
\item submeter ao orientador a especificação do Algoritmo 2, que entrará na
  seção~\ref{sec:classificacao};
\item implementar e executar o Algoritmo 2 sobre o corpus de resumos;
\item comparar o desempenho com o do \textit{baseline}, em pares por
  repetição, e analisá-lo por recorte, por idioma e por classe, em particular
  na classe manutenção;
\item avaliar a sumarização guiada por domínio, com o candidato C como prompt e
  o prompt neutro como controle (seção~\ref{sec:prompt}).
\end{enumerate}

\begin{thebibliography}{99}

\bibitem[Almeida et al. 2025]{almeida2025}
Almeida, T. S., Pires, R., Abonizio, H., Nogueira, R. e Pedrini, H. (2025).
PoETa v2: Toward More Robust Evaluation of Large Language Models in
Portuguese. \textit{arXiv:2511.17808}.

\bibitem[Del Piccolo 2026]{delpiccolo2026}
Del Piccolo, G. W. (2026). \textit{selic-direction-prediction-summaries}.
Repositório GitHub.
\url{https://github.com/giovanniwillik/selic-direction-prediction-summaries}.
Acesso em: 13 nov. 2026.

\bibitem[Devlin et al. 2019]{devlin2019}
Devlin, J., Chang, M.-W., Lee, K. e Toutanova, K. (2019). BERT: Pre-training of
Deep Bidirectional Transformers for Language Understanding. Em
\textit{Proceedings of NAACL-HLT 2019}, pp. 4171--4186.

\bibitem[Fasolo et al. 2022]{fasolo2022}
Fasolo, A. M., Graminho, F. M. e Bastos, S. B. (2022). Seeing the forest for the
trees: using hLDA models to evaluate communication in Banco Central do Brasil.
\textit{BIS Working Papers}, n.~1021. Bank for International Settlements.

\bibitem[Forti 2025]{forti2025}
Forti, R. A. (2025). \textit{Predição da direção futura da taxa SELIC a partir
de atas de reunião do COPOM}. Trabalho de Conclusão de Curso, EACH/USP. Código e
dados: \url{https://github.com/RyanForti/Selic-Direction-Prediciton}. Acesso em:
13 nov. 2026.

\bibitem[Goldberg 1991]{goldberg1991}
Goldberg, D. (1991). What Every Computer Scientist Should Know About
Floating-Point Arithmetic. \textit{ACM Computing Surveys}, 23(1), pp. 5--48.

\bibitem[Grattafiori et al. 2024]{grattafiori2024}
Grattafiori, A. et al. (2024). The Llama 3 Herd of Models.
\textit{arXiv:2407.21783}.

\bibitem[Holtzman et al. 2020]{holtzman2020}
Holtzman, A., Buys, J., Du, L., Forbes, M. e Choi, Y. (2020). The Curious Case
of Neural Text Degeneration. Em \textit{ICLR 2020}.

\bibitem[Jaccard 1912]{jaccard1912}
Jaccard, P. (1912). The distribution of the flora in the alpine zone.
\textit{New Phytologist}, 11(2), pp. 37--50.

\bibitem[Jingyi e Liu 2025]{jingyi2025}
Jingyi, F. X. e Liu, L. (2025). Can We Reliably Predict the Fed's Next Move? A
Multi-Modal Approach to U.S. Monetary Policy Forecasting.
\textit{arXiv:2506.22763}.

\bibitem[Kim et al. 2024]{kim2024}
Kim, Y., Chang, Y., Karpinska, M., Garimella, A., Manjunatha, V., Lo, K.,
Goyal, T. e Iyyer, M. (2024). FABLES: Evaluating faithfulness and content
selection in book-length summarization. Em \textit{1st Conference on Language
Modeling (COLM 2024)}. arXiv:2404.01261.

\bibitem[L'Ecuyer et al. 2002]{lecuyer2002}
L'Ecuyer, P., Simard, R., Chen, E. J. e Kelton, W. D. (2002). An
object-oriented random-number package with many long streams and substreams.
\textit{Operations Research}, 50(6), pp. 1073--1075.

\bibitem[Lin 2004]{lin2004}
Lin, C.-Y. (2004). ROUGE: A Package for Automatic Evaluation of Summaries. Em
\textit{Proceedings of the ACL Workshop on Text Summarization Branches Out},
pp. 74--81.

\bibitem[Mutasodirin e Prasojo 2024]{mutasodirin2024}
Mutasodirin, M. A. e Prasojo, R. E. (2024). Investigating Text Shortening
Strategy in BERT: Truncation vs Summarization. \textit{arXiv:2403.12799}.

\bibitem[Pappagari et al. 2019]{pappagari2019}
Pappagari, R., \.{Z}elasko, P., Villalba, J., Carmiel, Y. e Dehak, N. (2019).
Hierarchical Transformers for Long Document Classification. Em
\textit{Automatic Speech Recognition and Understanding Workshop, 2019}.
arXiv:1910.10781.

\bibitem[Rehman et al. 2025]{rehman2025}
Rehman, T., Ghosh, S., Das, K., Bhattacharjee, S., Sanyal, D. K. e
Chattopadhyay, S. (2025). Evaluating LLMs and Pre-trained Models for Text
Summarization Across Diverse Datasets. \textit{arXiv:2502.19339}.

\bibitem[Reimers e Gurevych 2019]{reimers2019}
Reimers, N. e Gurevych, I. (2019). Sentence-BERT: Sentence Embeddings using
Siamese BERT-Networks. Em \textit{Proceedings of EMNLP-IJCNLP 2019},
pp. 3982--3992.

\bibitem[Shah et al. 2023]{shah2023}
Shah, A., Paturi, S. e Chava, S. (2023). Trillion Dollar Words: A New Financial
Dataset, Task \& Market Analysis. Em \textit{Proceedings of the 61st Annual
Meeting of the Association for Computational Linguistics (ACL 2023)}.
arXiv:2305.07972.

\bibitem[Souza et al. 2020]{souza2020}
Souza, F., Nogueira, R. e Lotufo, R. (2020). BERTimbau: Pretrained BERT Models
for Brazilian Portuguese. Em \textit{Proceedings of BRACIS 2020}.

\bibitem[Sun et al. 2019]{sun2019}
Sun, C., Qiu, X., Xu, Y. e Huang, X. (2019). How to Fine-Tune BERT for Text
Classification? Em \textit{Proceedings of CCL 2019}, pp. 194--206.

\end{thebibliography}

\end{document}
